%===================================================================================================================
\subsection{Zitieren}
Sobald auf fremde Literaturquellen referenziert wird, ist die Verwendung einer korrekten Zitierweise und der Aufbau eines Literatur- und Quellenverzeichnis erforderlich.
Die Inhalte fremder Quellen können in einem direkten oder indirekten Zitat wiedergegeben werden.

Bei der wörtlichen Übernahme einer fremden Quelle im eigenen Text wird von einem direkten Zitat gesprochen.
Dieses steht immer in Anführungszeichen und muss buchstaben- und zeichengetreu übernommen werden, auch wenn das Original Rechtschreib- oder sonstige Fehler enthält.

Ein indirektes Zitat liegt vor, wenn fremde Gedanken und Ausführungen sinngemäß übernommen werden.
Es steht nie in Anführungszeichen und wird z.\,B. mit einem "`vgl."' referenziert.

Man spricht von einem Vollbeleg, wenn ein gesondertes Literaturverzeichnis vorliegt und auf die Einträge im Literaturverzeichnis verwiesen wird wie das in dieser wissenschaftlichen Arbeit der Fall ist.
Bei einem Kurzbeleg existiert kein gesondertes Literaturverzeichnis.
Die Quellen der Zitate werden dabei innerhalb des Textes oder in Fußnoten aufgeführt \cite{Kipp2012}.

Sinngemäß zitierte Sachverhalte werden durch die Quellenangabe z.\,B. am Ende eines Teilsatzes, eines Satzes oder Absatzes gekennzeichnet \cite[S.~109~ff.]{Hering2009}.
Eine wissenschaftliche Arbeit stützt sich nicht nur auf Schriftwerke, sondern auch auf andere geschützte Werke der Literatur, Wissenschaft und Kunst.
Der Textteil beinhaltet neben den eigenen und zitierten Inhalten auch Abbildungen, Tabellen und Gleichungen, welche an entsprechenden Stellen in den Text eingebunden werden \cite[S.~196]{Balz2008}.
Die Theorie ist mit entsprechenden Zitaten, Formeln, Tabellen und Bildern aus Fachliteratur zu belegen.
Dabei sollen grundsätzlich Primärquellen den Sekundärquellen vorgezogen werden.
Primärquellen oder Originalquellen sind eigenständige wissenschaftliche Arbeiten, die einen konkreten Betrachtungsgegenstand haben und zum Zeitpunkt der Veröffentlichung neue Erkenntnisse über diesen liefern \cite[S.~81]{Balz2008}.

"`Schwierige und komplexe Inhalte brauchen Wahrnehmungshilfen [...]: Dazu zählen Überschriften, Auszeichnungen, Bilder, Symbole und Strukturbilder, Tabellen und Diagramme, Fotos und multimediale Elemente"' (\cite{Balz2008}, S. 28).
Wenn Formeln und Gleichungen verwendet werden, so werden sie nicht in den Text eingebettet sondern nehmen eine eigene Zeile ein. 
Es ist üblich, die Gleichungen zusätzlich mit Gleichungsnummern am rechten Rand des Satzspiegels zu kennzeichnen, damit man sich im Text auf die jeweilige Gleichung beziehen kann.
Im gesamten Text kann auf die Formeln Bezug genommen werden indem ihre Nummer genannt wird.
Dabei werden die Klammern weggelassen \cite[S.~133~ff.]{Hering2009}.
Nachfolgend wird ein Beispiel für das korrekte Einbringen einer Formel mit dem dazugehörigen Querverweis im Text dargestellt.

Nach Gleichung \ref{eq:Gleichung1} errechnet sich die Gesamtsystemleistung $P_{ges}$ aus dem Produkt von Spannung $U_f$, Strom $I_f$ und der Anzahl der LEDs $n$.

\begin{equation}
P_{ges} = P_E = I_f \cdot U_f \cdot n
\label{eq:Gleichung1}
\end{equation}

\noindent
Formeln treten selten als einzelne Formel wie in Gleichung \ref{eq:Gleichung1} auf.
Meist erscheinen sie im Rahmen umfangreicherer Berechnungen mit mehreren Formeln und Umformungen.
Innerhalb des Textes sollten die Formelgrößen erklärt, Formelumformungen beschrieben und die Ergebnisse aufgegriffen werden.
Die Einflussgrößen, die in einer Formel auftreten, werden im Text kurz vor oder direkt nach der Formel benannt und erläutert \cite[S.~134]{Hering2009}.
Es folgt ein Beispiel für eine umfangreichere Berechnung.

Die LEDs sind mit einem Durchflussstrom von $700\,mA$ und einer Durchflussspannung von $3,31V$ klassifiziert.
Um die benötigte Systemleistung $P_{ges}$ aufzubringen, wird von dem typischen Durchflussstrom von $700\,mA$ ausgegangen.
Bei diesem stellt sich nach Gleichung \ref{eq:Spannung319V} eine Durchflussspannung von $3,19\,V$ ein und erreicht bei der Anzahl $n$ von 14 LEDs somit die Gesamtsystemleistung von $31,25\,W$.
%
\begin{align}
P_{ges} = P_E &= I_f \cdot U_f \cdot n \\
\Leftrightarrow U_f	&= \frac{P_{ges}}{I_f \cdot n}\\
	  &= \frac{31,25W}{700mA \cdot 14}\\
	  &= 3,19V
\label{eq:Spannung319V}
\end{align}

Gleichungen sowie Zahlen mit Einheiten werden kursiv hervorgehoben, falls diese im Textteil erscheinen.
Mit Latex gesetzte Formeln sehen einheitlich aus, da Formelzeichen in der mathematischen Umgebung grundsätzlich kursiv dargestellt werden.

%===================================================================================================================
\subsection{Quellen}
Fertigen Sie Ihr Literaturverzeichnis nach DIN ISO 690:2013-10 ist an.
Die DIN ISO 690:2013-10 sieht im Literaturverzeichnis zwingend die Nennung des Urhebers, Sachtitels, Auflage und Erscheinungsvermerke wie Erscheinungsort, Verlag und Erscheinungsjahr einer Literaturquelle vor.
\newpage
Dies sind Beispiele für Zitate nach DIN ISO 690:2013-10:

\textbf{Bücher}
  \begin{tabbing}
    [1]x \= Mitte  \kill
    $[1]$ \> \textsc{Paul, R.}: \textit{Elektrotechnik 1.} Berlin: Springer-Verlag, 3. Auflage 1993.\\
		\> ISBN: 978-0-387557-53-3\\
		$[2]$ \> \textsc{Balzert, H. u.a.}: \textit{Wissenschaftliches Arbeiten.} \\
		\> Witten: W3L-Verlag, 2. Auflage 2008. ISBN 987-3-937137-59-9\\
		$[3]$ \> \textsc{Eichler, H.-J.; Kronfeldt, H.-D.; Sahm, J.}: \textit{Das neue Physikalische Grund-}\\
		\>\textit{praktikum.} Berlin, Heidelberg: Springer-Verlag, 3. Auflage 2016.\\
		\> ISBN: 978-3-662490-22-8
  \end{tabbing}

\textbf{Internetquellen}
  \begin{tabbing}
    [1]x \= Mitte  \kill
    [4] \> \textsc{Wikipedia}: \textit{Feldlinien.} URL: \url{https://de.wikipedia.org/wiki/Feldlinien}, \\
		\>Abruf: 12.02.2022\\
  \end{tabbing}

\textbf{Zeitschriften}
  \begin{tabbing}
    [1]x \= Mitte  \kill
    [5] \> \textsc{Zhong, X.;Lewis, J. M.}: \textit{A New Method for Autonomous Robot Calibration.} \\
		\>In: Proc. of the IEEE Int. Conf. on Robotics and Automation. \\
		\>IEEE, 1995, S. 1790-1795 Vol.2. DOI: 10.1109/ROBOT.1995.525529.\\
  \end{tabbing}



